% cSpell:locale en,pt-BR - Adicionado para suporte ao corretor ortográfico
% Extensão Code Spell Checker no VS Code (muito útil!)
\chapter{Por onde seguir}\label{cap:conclusao}

Se você leu até aqui antes de começar a escrever, o caminho é curto: copie o arquivo de partida do
seu tipo, compile-o sem alterar nada, confira a capa que saiu e só então comece.

Duas coisas valem ser conferidas antes da entrega, porque nenhuma das duas produz erro de
compilação:

\begin{itemize}
    \item o \textbf{tipo declarado} no preâmbulo, que governa a natureza e a folha de aprovação
          (\refcomp{}{cap:estrutura});
    \item o \textbf{arquivo de referências}, que precisa ser o seu. O \texttt{referencias.bib} deste
          repositório é fictício de propósito, e serve apenas aos exemplos deste manual
          (\refcomp{}{sec:bibliografia}).
\end{itemize}

\section{Quando algo não sai como devia}\label{sec:problemas}

Compilação limpa não prova saída correta. Antes de concluir que o modelo está errado, abra o PDF e
leia a página: boa parte dos enganos --- rótulo antes da legenda, tipo de documento trocado, curso
padrão não substituído --- compila sem qualquer aviso.

Problemas no pacote, dúvidas sobre a formatação exigida e sugestões de melhoria devem ser
registrados no repositório do modelo:

\begin{center}
    \url{https://github.com/WndDnz/modelo-relatorio-unifei-ict}
\end{center}

\noindent Ao relatar, informe o tipo de documento declarado e anexe, se possível, um exemplo mínimo
que reproduza o problema.
